% from ega1-2-fr.tex
\oldpage[I]{100}
correspondance biunivoque canonique entre les $f$-morphismes
$\mathscr{F}\to\mathscr{G}$ (\hyperref[0.4.4.1-fr]{0, 4.4.1}) et les
$A$-homomorphismes
$M\to(\Gamma(X,\mathscr{G}))_{[\varphi]}$.
\end{proposition}

En effet, en raisonnant comme dans (\hyperref[I.2.2.4-fr]{2.2.4}), on est
aussitôt ramené au cas où $X$ est affine et la proposition résulte alors de
(\hyperref[I.1.6.3-fr]{1.6.3}) et (\hyperref[I.1.3.8-fr]{1.3.8}).

\begin{env}[2.2.6]
\label{I.2.2.6-fr}
On dit qu'un morphisme de préschémas
$(\psi,\theta):(X,\mathscr{O}_X)\to(Y,\mathscr{O}_Y)$ est \emph{ouvert}
(resp. \emph{fermé}), si pour toute partie ouverte $U$ de $X$ (resp. toute
partie fermée $F$ de $X$), $\psi(U)$ est ouvert dans $Y$ (resp. $\psi(F)$
fermé dans $Y$). On dit que $(\psi,\theta)$ est \emph{dominant} si $\psi(X)$
est dense dans $Y$, \emph{surjectif} si $\psi$ est surjectif. On notera que
ces conditions ne font intervenir que l'application continue $\psi$.
\end{env}

\begin{proposition}[2.2.7]
\label{I.2.2.7-fr}
Soient
\[
  f=(\psi,\theta):(X,\mathscr{O}_X)\to(Y,\mathscr{O}_Y),
  \qquad
  g=(\psi',\theta'):(Y,\mathscr{O}_Y)\to(Z,\mathscr{O}_Z)
\]
deux morphismes de préschémas.
\begin{enumerate}
  \item[\rm(i)] Si $f$ et $g$ sont tous deux ouverts (resp. fermés, dominants,
    surjectifs), il en est de même de $g\circ f$.
  \item[\rm(ii)] Si $f$ est surjectif, et si $g\circ f$ est fermé, $g$ est
    fermé.
  \item[\rm(iii)] Si $g\circ f$ est surjectif, $g$ est surjectif.
\end{enumerate}
\end{proposition}

Les assertions (i) et (iii) sont évidentes. Posons
$g\circ f=(\psi'',\theta'')$. Si $F$ est fermé dans $Y$,
$\psi^{-1}(F)$ est fermé dans $X$, donc
$\psi''(\psi^{-1}(F))$ est fermé dans $Z$~; mais comme $\psi$ est surjectif,
$\psi(\psi^{-1}(F))=F$, donc
$\psi''(\psi^{-1}(F))=\psi'(F)$, ce qui démontre (ii).

\begin{proposition}[2.2.8]
\label{I.2.2.8-fr}
Soient $f=(\psi,\theta)$ un morphisme
$(X,\mathscr{O}_X)\to(Y,\mathscr{O}_Y)$, $(U_\alpha)$ un recouvrement ouvert
de $Y$. Pour que $f$ soit ouvert (resp. fermé, surjectif, dominant), il faut
et il suffit que sa restriction à chacun des préschémas induits
$(\psi^{-1}(U_\alpha),\mathscr{O}_X|\psi^{-1}(U_\alpha))$, considérée comme
un morphisme de ce préschéma induit dans le préschéma induit
$(U_\alpha,\mathscr{O}_Y|U_\alpha)$, soit ouvert (resp. fermé, surjectif,
dominant).
\end{proposition}

La proposition résulte aussitôt des définitions, tenant compte du fait qu'une
partie $F$ de $Y$ est fermée (resp. ouverte, dense) dans $Y$ si et seulement
si chacun des ensembles $F\cap U_\alpha$ est fermé (resp. ouvert, dense) dans
$U_\alpha$.

\begin{env}[2.2.9]
\label{I.2.2.9-fr}
Soient $(X,\mathscr{O}_X)$, $(Y,\mathscr{O}_Y)$ deux préschémas~; on suppose
que $X$ et $Y$ ont un même nombre \emph{fini} de composantes irréductibles
$X_i$ (resp. $Y_i$) $(1\leq i\leq n)$~; soit $\xi_i$ (resp. $\eta_i$) le
point générique de $X_i$ (resp. $Y_i$)
(\hyperref[I.2.1.5-fr]{2.1.5}). On dit qu'un morphisme
\[
  f=(\psi,\theta):(X,\mathscr{O}_X)\to(Y,\mathscr{O}_Y)
\]
est \emph{birationnel} si, pour tout $i$, $\psi^{-1}(\eta_i)=\{\xi_i\}$ et
$\theta_{\xi_i}^{\sharp}:\mathscr{O}_{\eta_i}\to\mathscr{O}_{\xi_i}$ est un
\emph{isomorphisme}. Il est clair qu'un morphisme birationnel est
\emph{dominant} (\hyperref[0.2.1.8-fr]{0, 2.1.8}), donc surjectif s'il est
aussi fermé.
\end{env}

\begin{notation}[2.2.10]
\label{I.2.2.10-fr}
Dans toute la suite de cet ouvrage et lorsque cela ne risquera pas de créer
des confusions, nous \emph{supprimerons} dans la notation d'un préschéma
(resp. d'un morphisme) le faisceau structural (resp. le morphisme de faisceaux
structuraux). Si $U$ est une partie ouverte de l'espace sous-jacent d'un
préschéma $X$, lorsqu'on parlera de $U$ comme d'un préschéma, il s'agira
toujours du préschéma induit sur $U$.
\end{notation}

\subsection{Recollement de préschémas}
\label{subsection:I.2.3-fr}

\begin{env}[2.3.1]
\label{I.2.3.1-fr}
